\subsection{Setting}
\paragraph{Purpose}
In supervised learning the correct output is given for each input. In reinforcement learning
(RL) an \tB{agent} interacts with an \tB{environment}: at each step it observes a state,
chooses an action, and receives a scalar \tB{reward}. No one tells it the right action, it has
to discover which \uB{sequences of actions} give the most reward, and a reward may come long
after the action that caused it (\emph{credit assignment}).

\paragraph{Markov Decision Process}
An MDP is a tuple $\left(\mathcal{S},\mathcal{A},P,R,\gamma\right)$:
\begin{itemize}
    \item $\mathcal{S}$ the states and $\mathcal{A}$ the actions;
    \item $P(s'\,|\,s,a)=\prob{S_{t+1}=s'|S_{t}=s,A_{t}=a}$ the transition probabilities,
        with the Markov assumption: the future depends on the past only through
        $(s,a)$;
    \item $R(s,a)$ the expected reward;
    \item $\gamma\in[0,1[$ the discount factor.
\end{itemize}
It is a \tB{Markov chain whose transitions are controlled by the actions} (see Markov
models), and with a single action it reduces to a Markov chain.

\paragraph{Policy and return}
A \textbf{policy} $\pi(a\,|\,s)$ is the probability of choosing the action $a$ in the state
$s$. Following $\pi$ generates a trajectory $\tau=(s_{0},a_{0},r_{1},s_{1},a_{1},\dots)$.
The goal is to maximize the expected discounted \textbf{return}
\begin{center}
    \frB{$G_{t}=\su{k=0}{\infty}\gamma^{k}\,r_{t+k+1},\qquad
    J(\pi)=\mathbb{E}_{\tau\sim\pi}\left[G_{0}\right]$}
\end{center}
$\gamma<1$ makes the sum finite and favors immediate rewards. Note that $G_{t}=r_{t+1}+
\gamma G_{t+1}$.

\subsection{Value Functions and Bellman Equations}
\paragraph{Definitions}
\begin{center}
    $V^{\pi}(s)=\mathbb{E}_{\pi}\left[G_{t}\,|\,S_{t}=s\right],\qquad
    Q^{\pi}(s,a)=\mathbb{E}_{\pi}\left[G_{t}\,|\,S_{t}=s,A_{t}=a\right]$
\end{center}
The state value $V^{\pi}$ is the expected return from $s$, the action value $Q^{\pi}$ is the
expected return when taking $a$ first and then following $\pi$. They are linked by
$V^{\pi}(s)=\su{a}{}\pi(a|s)Q^{\pi}(s,a)$, and the \textbf{advantage}
$A^{\pi}(s,a)=Q^{\pi}(s,a)-V^{\pi}(s)$ tells how much better $a$ is than the average action
of $\pi$.

\paragraph{Bellman expectation equation} Using $G_{t}=r_{t+1}+\gamma G_{t+1}$:
\begin{center}
    \frB{$V^{\pi}(s)=\su{a}{}\pi(a|s)\left[R(s,a)+\gamma\su{s'}{}P(s'|s,a)V^{\pi}(s')\right]$}
\end{center}
The value of a state is the immediate reward plus the discounted value of the next state:
a \uB{recursive} definition, a linear system for a fixed $\pi$.

\paragraph{Bellman optimality equation} The optimal value functions $V^{*},Q^{*}$ satisfy
\begin{center}
    \frB{$Q^{*}(s,a)=R(s,a)+\gamma\su{s'}{}P(s'|s,a)\max_{a'}Q^{*}(s',a')$},\qquad
    $V^{*}(s)=\max_{a}Q^{*}(s,a)$
\end{center}
and the optimal policy is greedy with respect to $Q^{*}$: $\pi^{*}(s)=\argmax_{a}Q^{*}(s,a)$.
The $\max$ makes the equation non-linear.

\paragraph{Dynamic programming} If the model $(P,R)$ is known, \emph{value iteration}
repeats the update below until convergence:
\begin{center}
    $V(s)\leftarrow\max_{a}\left[R(s,a)+\gamma\su{s'}{}P(s'|s,a)V(s')\right]$
\end{center}
In RL the model is usually \tB{unknown} and the agent has to learn from sampled
interactions, which is what the next methods do.

\paragraph{Exploration versus exploitation} To learn, the agent must try actions that look
suboptimal (\emph{explore}), but to collect reward it must use what it knows
(\emph{exploit}). The simplest rule is $\varepsilon$-greedy: take a random action with
probability $\varepsilon$, the greedy one otherwise.
